\section{Related Work}
\label{sec:related-work}
%% ====================================================================

We organize the related work into ontologies for machine learning, ontologies
for materials science, and knowledge graphs and
FAIR~\cite{wilkinson2016fair,Wilkinson2019appendum} platforms in the domain.

\subsection{Ontologies for Machine Learning}

ML-Schema~\cite{publio2018mlschema}, from a W3C community group, provides
top-level classes for ML algorithms, datasets, and experiments.
DMOP~\cite{keet2015dmop} supports ontology-driven algorithm selection and
hyperparameter optimization for data mining workflows. The MEX
vocabulary~\cite{esteves2015mex} provides a lightweight RDF vocabulary for ML
experiment metadata, built on PROV-O. The Common Metadata Ontology
(CMO)~\cite{foltin2024cmo} integrates pipeline metadata from multiple ML
platforms. These efforts address ML workflows in general but do not cover the
domain-specific concepts of MLIP research: algorithm families with
physics-specific hyperparameters (cutoff radii, angular/radial basis functions),
training set construction strategies, and DFT reference calculation settings.

\subsection{Ontologies for Materials Science}

The European Materials Modelling Ontology (EMMO)~\cite{horsch2020emmo} provides
a top-level ontology grounded in mereocausality theory. The Materials Design
Ontology (MDO)~\cite{li2024mdo} is an OWL\,2~DL ontology with four modules
(Core, Structure, Calculation, Provenance); it has been used for
SPARQL-based integrated querying over the Materials
Project~\cite{jain2013mp} and OQMD~\cite{saal2013oqmd}. PMDco~\cite{bayerlein2024pmdco}
provides a mid-level ontology aligned with BFO. CMSO and
ASMO~\cite{menon2024atomrdf} describe atomic structures and simulation methods.
A recent survey by Norouzi et al.~\cite{norouzi2024survey} analyzes
60~ontologies in the field.

These ontologies model the physical domain at various levels of granularity but
do not address the specific vocabulary needed for MLIP training: which
hyperparameters govern each algorithm family, how training datasets are
constructed and what DFT settings were used, and how to structure benchmark
comparisons across algorithms, materials, and properties.

\subsection{Metadata Schemas and Formats beyond the Semantic Web}

Several major platforms in computational materials science define metadata
schemas for organizing simulation data, but express them in
platform-specific formats rather than Semantic Web standards.
NOMAD~\cite{draxl2023nomad} has developed the NOMAD Metainfo, a rich
metadata schema that normalizes outputs from diverse simulation codes; it is
evolving toward alignment with EMMO but is not yet published as an OWL
ontology with a SPARQL endpoint.
The Materials Project~\cite{jain2013mp},
AFLOW~\cite{curtarolo2012aflow}, and OQMD~\cite{saal2013oqmd} each expose
data through custom REST APIs with their own JSON schemas.
OPTIMADE~\cite{andersen2021optimade} standardizes API access across these
databases through a JSON:API specification, enabling cross-database queries
but without the formal semantics, reasoning capabilities, or federation
mechanisms that the Resource Description Framework
(RDF)~\cite{rdf11concepts} and OWL~\cite{owl2overview} provide.
AiiDA~\cite{huber2020aiida} automatically tracks full data provenance as
a directed acyclic graph and, together with Materials
Cloud~\cite{talirz2020materialscloud}, provides infrastructure for data
dissemination; however, the provenance model uses AiiDA's own graph
representation rather than PROV-O or other Semantic Web standards.
Kadi4Mat~\cite{brandt2021kadi4mat} is a virtual research environment for
FAIR research data management, providing hierarchical collections with
group-based access control and metadata records; it serves, for example, as
the file repository in the Circular Factory
CRC~1574~\cite{thapa2025inf}. While effective for file and metadata
management, it does not provide SPARQL endpoints or ontology-driven data
editing.

The existence of these schemas demonstrates that the materials science
community recognizes the need for structured metadata. However, because each
platform defines its own format, federated querying across databases and
formal reasoning over the combined metadata are not possible. The
\MLIPsOntology{} addresses this gap by providing a shared vocabulary
expressed in OWL and queryable via SPARQL, designed to complement these
platforms by giving explicit, machine-interpretable semantics to the metadata
they already capture in ad-hoc formats. Recent work by Azoc{\'a}r Guzm{\'a}n
et al.~\cite{azocarguzman2026workflows} confirms this direction,
demonstrating FAIR-compliant RDF workflows using CMSO, ASMO, and PROV-O
within pyiron---showing that atomistic simulation communities are
increasingly adopting Semantic Web standards.

%% ====================================================================
